\begin{center}
\begin{minipage}{0.98\linewidth}
\begin{tcolorbox}[colback=white,colframe=CStatement!25,boxrule=0.4pt,
  arc=1mm,left=3mm,right=3mm,top=3mm,bottom=3mm]
\centering
\textbf{点到直线的距离，取垂线段 PH}\par\medskip
\begin{tikzpicture}[font=\small,>=Stealth]
\begin{scope}[x=1.7cm,y=1.7cm]
\coordinate (P) at (0,3);\coordinate (H) at (1.44,1.08);
\draw[CStatement,thick] (-.8,-.6)--(3.8,2.85) node[above left] {$l:Ax+By+C=0$};
\draw[CExample,very thick] (P)--(H);
\draw[CProof] (1.296,1.272)--(1.488,1.416)--(1.632,1.224);
\fill[CExample] (P) circle (2pt);\fill[CProof] (H) circle (2pt);
\node[above left] at (P) {$P(x_0,y_0)$};
\draw[gray] (H)--(2.05,1.15);\node[right] at (2.05,1.15) {$H$};
\node[above right,text=CExample] at (.72,2.04) {$d=PH$};
\end{scope}
\end{tikzpicture}\par\medskip
$H$ 在直线 $l$ 上，$PH\perp l$，距离就是 $PH$ 的长度。若 $P$ 本来在 $l$ 上，则距离为 $0$。
\end{tcolorbox}
\end{minipage}
\end{center}
